\subsection{Type Safety}
\label{sec:type}

Move's nominal typing means two structs with identical layout but different names are distinct types.
A \emph{type-confusion} exploit breaks that distinction, and is most damaging when one of the two types is privileged.
Consider the following example:

\begin{MoveBox}
module 0x1::admin {
  struct Capability(address);
  public fun admin_action(c: &Capability) {
    /* gated by Capability */
  }
}

module 0x2::attacker {
  struct FakedCapability(address);
}
\end{MoveBox}

The attacker freely constructs a |FakedCapability| (it owns the type) and only needs a route that delivers it to |admin_action| as if it were a |Capability|.

\subsubsection{Example} Back when Move was developed at Meta/Facebook, vector operations were added to the VM.
The developer forgot to add a type check in the bytecode verifier for the vector element type on push.
The compiler did perform this check, so no test case via the source language could reach the problem.
However, it was possible to manually construct bytecode, skipping the compiler, which did the following (conceptually):

\begin{MoveBox}
let v: vector<Capability> = vector::empty();
v.push_back(FakedCapability(victim)); // should be rejected
admin_action(&v[0]);  // runs on a forged token
\end{MoveBox}

The two types share a layout, so every downstream read succeeds, and the gated action runs on a token its owner never minted. This bug was found as part of the Aptos bug bounty program before the network went to mainnet.


\subsubsection{Performing the Checks}
In order to prevent type confusion, we perform type checks at runtime.
To this end a type stack is maintained in parallel to the operand stack of the Move VM.
Whenever a value is pushed on the operand stack, its type is also pushed on the type stack.
This type is known since the Move VM has typed registers.
Whenever a value is popped from the stack and stored in a register or consumed by an operation, the corresponding type is popped as well and compared with the expected type, creating a runtime error if type confusion is detected.


%$\theta ::= \epsilon \mid \theta \cdot \tau$ is added to the execution environment in Fig.~\ref{fig:semantics}.
%Let $r$ be a reference such that $r = (\textrm{index}(F), l)$ on $\COPYLOC l$ and $\MOVELOC l$, and the reference %on the stack for $\READREF$, then $\tau(r)$ is pushed on the type stack.
% When $\WRITEREF$ is executed, we pop the type $\tau'$ from the type stack and check $\tau' = \tau(r)$; if the check fails, type confusion has been detected.





%%% Local Variables:
%%% mode: latex
%%% TeX-master: "main"
%%% End:
